\documentclass[11pt]{article}
\usepackage[margin=1in]{geometry}
\usepackage{amsmath,amssymb,amsthm}
\usepackage{fontspec}
\setmonofont{DejaVu Sans Mono}[Scale=0.78]
\newfontfamily\ellfont{FreeMono}[Scale=0.85]
\usepackage{newunicodechar}
\newunicodechar{ℓ}{{\ellfont ℓ}}
\usepackage{fvextra}
\usepackage{microtype}
\usepackage{enumitem}
\setlist{nosep,leftmargin=*}
\usepackage{xurl}
\usepackage[colorlinks=true,allcolors=blue]{hyperref}
\sloppy
\newtheorem{theorem}{Theorem}
\newtheorem{corollary}[theorem]{Corollary}
\theoremstyle{remark}
\newtheorem{remark}[theorem]{Remark}
\newcommand{\Res}{\mathrm{Res}(\oplus)}
\newcommand{\BPHP}{\mathrm{BPHP}}
\newcommand{\Col}{\mathrm{Col}}
\setlength{\parskip}{0.3em}
\title{Bit PHP with more pigeons, and bit-colouring, in resolution over parities:\\
a corollary of Braun's lower bound, with a Lean proof}
\author{Prespatial Labs\thanks{Draft of 19 October 2026, prepared with Claude
(Anthropic). The mathematics is an immediate corollary of~\cite{Br26}; the Lean
proofs are checked by Lean's kernel against the formalization of~\cite{Br26}.
No human expert has reviewed this note.}}
\date{}
\begin{document}
\maketitle

\begin{abstract}
Braun~\cite{Br26} proved that every DAG-like $\Res$ refutation of the bit
pigeonhole principle with $n+1$ pigeons and $n=2^\ell$ holes has more than
$\exp(n/(32768\,\ell^2))$ lines. The lemmas of his formalization take the number
of pigeons as a parameter, and only the final calibration fixes it to $n+1$.
Calibrating for $m$ pigeons gives, for $\ell\ge32$, a bound of
$\exp(n^2/(8192\,m\,\ell^2))$ for every $m$ with $n<m\le n^2/(8192\,\ell^2(2\ell+5))$, superpolynomial while
$m=o(n^2/\log^3 n)$. A proper colouring turns the bit encoding of colouring a
graph into bit PHP by a substitution, so the same bound holds for every graph
whose chromatic number lies in that range. Both statements are proved in one
Lean file on top of the formalization of~\cite{Br26}, using only the standard
axioms.
\end{abstract}

\section{Statements}

$\Res$ works with linear clauses, disjunctions of affine equations over
$\mathbb F_2$. A refutation is a sequence of lines, each an initial clause or
derived from earlier lines, ending in the empty clause; its size is the number
of lines. We use the proof system of the formalization of~\cite{Br26}: initial
clauses, semantic weakening, resolution on any affine form, and any sound
inference from two earlier lines. It contains every standard rule set, so
lower bounds for it hold for standard DAG-like $\Res$.

Let $n=2^\ell$. The bit pigeonhole principle $\BPHP^m_n$ has variables
$b_{i,t}$ for $i\in[m]$ and $t\in[\ell]$. For every two pigeons $i<i'$ and
every label $z\in\{0,1\}^\ell$, it has the clause
$\bigvee_t (b_{i,t}\ne z_t)\vee\bigvee_t (b_{i',t}\ne z_t)$. For a graph
$G=(V,E)$, the bit encoding $\Col_\ell(G)$ of colouring $G$ with $2^\ell$
colours has variables $b_{v,t}$ for $v\in V$ and the same clause for every
edge $\{u,w\}$ and every label $z$. It is satisfiable exactly when $G$ is
$2^\ell$-colourable.

\begin{theorem}[more pigeons]\label{thm:more}
Let $\ell\ge 32$ and $n=2^\ell$, and let $m$ satisfy $n<m$ and
$8192\,m\,\ell^2(2\ell+5)\le n^2$. Every refutation of $\BPHP^m_n$ has more
than $\exp\bigl(n^2/(8192\,m\,\ell^2)\bigr)$ lines.
\end{theorem}

\begin{theorem}[transfer]\label{thm:transfer}
Let $f\colon V\to[m]$ be a proper colouring of $G$. The substitution
$b_{v,t}\mapsto b_{f(v),t}$ maps every refutation of $\Col_\ell(G)$, line by
line, to a refutation of $\BPHP^m_{2^\ell}$ with the same number of lines.
\end{theorem}

\begin{corollary}\label{cor:colouring}
If $\ell\ge32$, $2^\ell<\chi(G)$ and $8192\,\chi(G)\,\ell^2(2\ell+5)\le 4^\ell$,
every refutation of $\Col_\ell(G)$ has more than
$\exp\bigl(4^\ell/(8192\,\chi(G)\,\ell^2)\bigr)$ lines.
\end{corollary}

\begin{remark}[ranges and instances]\leavevmode
\begin{itemize}
\item At $\ell=32$ the range is $2^{32}<m\le 31{,}869{,}902{,}254\approx 7.42\cdot
2^{32}$. In general its upper end is about $n^2/(16384\,\ell^3)$.
\item The pigeon count $m=2n$ is in range for every $\ell\ge32$, with bound
$\exp(n/(16384\,\ell^2))$. The count $m=\lfloor n^{3/2}\rfloor$ is in range
for every $\ell\ge65$, with bound at least $\exp(\sqrt n/(8192\,\ell^2))$.
\item For $m=n+1$ the bound is $\exp(n^2/(8192(n+1)\ell^2))$, of the same form
as~\cite{Br26}. The constants were not optimized.
\item The Mycielskian of $K_n$ has chromatic number $n+1$ but no clique on
$n+1$ vertices~\cite{My55}. By Corollary~\ref{cor:colouring} its bit-colouring
with $n=2^\ell$ colours is still exponentially hard, although it contains no
copy of $\BPHP^{n+1}_n$ as a sub-formula.
\item The encoding assumes the number of colours is a power of two. For other
numbers, the clauses that exclude unused labels are not treated here.
\end{itemize}
\end{remark}

\section{Proofs}

\begin{proof}[Proof of Theorem~\ref{thm:more}]
This is the proof of \texttt{bitPHP\_exponential} in~\cite{Br26} with $m$ in
place of $n+1$, using two of its lemmas.
\begin{itemize}
\item Transfer (\texttt{usual\_bitPHP\_PC\_transfer}): a refutation with $S$
lines yields a polynomial-calculus derivation of $1$ of degree $D=12\ell+1$,
over at most $M=3S+\binom m2$ blocks of extension variables.
\item Exclusion (\texttt{affine\_family\_exclusion\_of\_square}): no such
derivation exists if $1\le k\le m$, $m\ln(4M)\le k^2$, $4(k-1)<n$ and
$2k(D+1)-1\le n$.
\end{itemize}
Only the first two conditions involve $m$. Take $k=\lfloor n/(32\ell)\rfloor$,
as in~\cite{Br26}; then $k\le n<m$ and $n\le 64\,\ell\,k$, and the last two
conditions hold.

Suppose $S\le e^x$ with $x=n^2/(8192\,m\,\ell^2)$. The hypothesis gives
$x\ge2\ell+5$, so $e^x\ge 2^{2\ell+5}=32n^2$. It also gives $m\le n^2$. Hence
$4M\le 12e^x+4m^2\le e^{2x}$, and
\[
 m\ln(4M)\le 2mx=\frac{n^2}{4096\,\ell^2}\le k^2 .
\]
The exclusion lemma now contradicts the derivation, so $S>e^x$.
\end{proof}

\begin{proof}[Proof of Theorem~\ref{thm:transfer}]
Map each affine form $g$ over the vertex variables to $y\mapsto g(y\circ\hat
f)$, where $\hat f(v,t)=(f(v),t)$, and map each line to the set of images of
its forms.
\begin{itemize}
\item The image of a line holds at $y$ exactly when the line holds at
$y\circ\hat f$. So the image of every inference is again sound: weakening and
two-premise steps map to the same steps, and resolution maps to a two-premise
step.
\item The initial clause for an edge $\{u,w\}$ and a label $z$ maps to the
clause of $\BPHP^m_{2^\ell}$ for the pigeons $f(u)$ and $f(w)$ and the label
$z$. This uses $f(u)\ne f(w)$; the clause is symmetric in the two pigeons.
\item The empty clause maps to the empty clause.\qedhere
\end{itemize}
\end{proof}

Corollary~\ref{cor:colouring} follows by taking $f$ an optimal colouring and
$m=\chi(G)$. With $m=n+1$, Theorem~\ref{thm:transfer} together with Theorem~1.1
of~\cite{Br26} already gives Braun's bound for every graph with chromatic
number $n+1$.

\section{Earlier work}

\begin{itemize}
\item \cite{Br26} proves the case $m=n+1$, with no restriction on the
refutation. Its moment construction needs no assumption on the number of
pigeons.
\item Byramji and Impagliazzo~\cite{BI25} bound bounded-depth $\Res$ for bit
PHP. A refutation of size $S$ and depth $D$ satisfies:
\begin{itemize}
\item for every $m>n$, $D\sqrt{\log S}=\Omega(n^{1.5})$ \cite[Theorem~1.1]{BI25};
\item for $m=n+k$ with $k\le\delta n$,
$\log S=\Omega(\min\{n^2/D,\,n^4/(kD^2)\})$ \cite[Theorem~1.3]{BI25}.
\end{itemize}
These bounds weaken as the depth grows. When the depth is not restricted, they
give no superpolynomial bound.
\item Efremenko, Garl\'ik and Itsykson~\cite{EGI} prove:
\begin{itemize}
\item a size bound for regular $\Res$ refutations of $\BPHP^{n+1}_n$;
\item a size bound of $2^{\Omega(n)}$ for tree-like refutations of weak bit
PHP;
\item a width bound for DAG-like refutations of weak bit PHP.
\end{itemize}
\end{itemize}
A brief search on 19 October 2026 found no superpolynomial size bound without
depth restriction for DAG-like $\Res$ refutations of bit PHP with more than
$n+1$ pigeons. It also found none for bit-colouring beyond complete graphs. The
statements here follow directly from the development of~\cite{Br26} and may
well be known to experts.

\section{Formal verification}
\begingroup\raggedright

The file \texttt{MorePigeonsColouring.lean}, in the namespace
\texttt{MorePigeonsColouring}, imports \texttt{claims.BitPHPExponential}. It
proves eight declarations:
\begin{itemize}
\item the main results:
\begin{itemize}
\item \texttt{bitPHP\_more\_pigeons} (Theorem~\ref{thm:more});
\item \texttt{colouring\_refutation\_transfer} (Theorem~\ref{thm:transfer});
\item \texttt{bitColouring\_more\_pigeons} (Corollary~\ref{cor:colouring});
\end{itemize}
\item witnesses against vacuity: \texttt{range\_nonempty},
\texttt{usualBitPHP\_unsat\_more}, \texttt{bitColouring\_unsat},
\texttt{bitColouring\_sat} and \texttt{complete\_graph\_is\_bitPHP}.
\end{itemize}
Each depends only on \texttt{propext}, \texttt{Classical.choice} and
\texttt{Quot.sound}.

A graph is a symmetric loopless relation on \texttt{Fin v}. Mathlib's
\texttt{SimpleGraph} is not part of the formalization's build, so it is not
used. The colouring clauses reuse \texttt{usualCNFClause u w z} for edges
$u<w$.

To replay:
\begin{enumerate}
\item Take a checkout of the \texttt{math-research} repository at commit
\texttt{8904bf09}.
\item Build it, for example with \texttt{lake build claims.BitPHPExponential}.
\item Copy the file to \texttt{formalization/claims/}.
\item From \texttt{formalization/}, run \texttt{lake build
claims.MorePigeonsColouring}.
\end{enumerate}
Every formalization file the drop-in depends on, the toolchain and the Lake
manifest are byte-identical at commit \texttt{25239211}, so the same holds
there.

Our checker also verified four controls, each of which the kernel rejects:
\begin{itemize}
\item lowering the constant $8192$ to $2048$ throughout the proof of
Theorem~\ref{thm:more}, which fails exactly at the step that calibrates $k$;
\item the transfer without properness;
\item unsatisfiability of bit PHP with as many pigeons as labels;
\item unsatisfiability of bit-colouring without a monochromatic edge.
\end{itemize}
The main formal statements follow.
\par\endgroup

\begin{footnotesize}
\VerbatimInput[breaklines=true,breakanywhere=true]{lean_statements.txt}
\end{footnotesize}

\begin{thebibliography}{EGI25}\raggedright
\bibitem[Br26]{Br26} K.~Braun, \emph{An exponential lower bound for the bit
pigeonhole principle in resolution over parities}, arXiv:2609.23015v1, 2026.
Formalization: \url{https://github.com/kbr-/math-research}, commit
\texttt{8904bf09}.
\bibitem[BI25]{BI25} F.~Byramji and R.~Impagliazzo, \emph{Lower bounds for bit
pigeonhole principles in bounded-depth resolution over parities},
arXiv:2511.20023v1, 2025; earlier version ECCC TR25-118.
\bibitem[EGI25]{EGI} K.~Efremenko, M.~Garl\'ik and D.~Itsykson, \emph{Lower
bounds for regular resolution over parities}, SIAM J.\ Comput., 2025; ECCC
TR23-187.
\bibitem[My55]{My55} J.~Mycielski, \emph{Sur le coloriage des graphes},
Colloq.\ Math.\ 3 (1955), 161--162.
\end{thebibliography}
\end{document}
